\section*{अध्याय \ref{ch:b3:dna-repair} --- जीनोम स्थायित्व: DNA क्षति, मरम्मत और पुनर्संयोजन}

\begin{solution}{exo:b3:dna-repair:1}
यूरैसिल: स्वतःस्फूर्त (C का \omterm{def:b3:dna-repair:lesions}{विअमीनन}), यूरैसिल-DNA \omterm{def:b3:dna-repair:excision}{ग्लाइकोसिलेज़} से
\omterm{def:b3:dna-repair:excision}{क्षार-उच्छेदन मरम्मत}। थाइमीन द्विलक: परिवेशी (UV),
\omterm{def:b3:dna-repair:excision}{न्यूक्लियोटाइड-उच्छेदन मरम्मत}। \omterm{def:b3:dna-repair:lesions}{8-ऑक्सोग्वानीन}: स्वतःस्फूर्त
(ऑक्सीकरण), 8-ऑक्सोG \omterm{def:b3:dna-repair:excision}{ग्लाइकोसिलेज़} से \omterm{def:b3:dna-repair:excision}{क्षार-उच्छेदन मरम्मत}। द्विशाख के
पीछे G--T असंगति: प्रतिकृति त्रुटि, \omterm{prop:b3:dna-repair:fidelity}{असंगति मरम्मत}। G1 में \omterm{def:b3:dna-repair:lesions}{द्विलड़ी
विखंडन}: परिवेशी (विकिरण) या स्वतःस्फूर्त; \omterm{def:b3:dna-repair:dsb}{असमजात सिरा-संयोजन}, क्योंकि
कोई सहोदरी क्रोमैटिड उपलब्ध नहीं।
\end{solution}

\begin{solution}{exo:b3:dna-repair:2}
न्यूक्लियोटाइड-उच्छेदन उस भौतिक गुण को पहचानता है जो सारी स्थूल
\omterm{def:b3:dna-repair:lesions}{विक्षतियाँ} साझा करती हैं --- कुंडली का विरूपण, या कोई रुकी RNA
पॉलीमरेज़ --- और ऐसा चकत्ता हटा देता है जिसमें उसका कारण भी आ जाता है।
छोटी \omterm{def:b3:dna-repair:lesions}{विक्षतियाँ} (यूरैसिल, 8-ऑक्सोG, क्षारकृत क्षार) कुंडली को शायद ही
विरूपित करती हैं, इसलिए उन्हें उनकी रसायन से पहचानना पड़ता है, हर
प्रकार के गलत क्षार के लिए एक \omterm{def:b3:dna-repair:excision}{ग्लाइकोसिलेज़}।
\end{solution}

\begin{solution}{exo:b3:dna-repair:3}
किसी असंगति में दो सामान्य क्षार होते हैं, और गलत केवल वही है जो नई
लड़ी पर है; तंत्र को पता होना चाहिए कि कौन-सी लड़ी नई है।
\emph{E.~coli} जनक-लड़ी को GATC पर मेथिल समूहों से चिह्नित करता है,
जो नई लड़ी पर कुछ मिनटों तक नहीं होते; MutH अमेथिलित लड़ी में कटान
लगाती है। किसी \emph{dam} उत्परिवर्ती में कोई भी लड़ी मेथिलित नहीं
होती: MutH किसी में भी कटान लगाती है, मरम्मत सही क्षार को उतनी ही बार
हटाती है जितनी गलत को --- यानी आधी त्रुटियाँ उत्परिवर्तन के रूप में
जमा देती है (एक उत्परिवर्तक) --- और पास-पास के GATC स्थलों पर दोनों
लड़ियों में लगे कटान \omterm{def:b3:dna-repair:lesions}{द्विलड़ी विखंडन} बना देते हैं।
\end{solution}

\begin{solution}{exo:b3:dna-repair:4}
लगभग $35\times 2 = 70$ \omterm{def:b3:dna-repair:lesions}{द्विलड़ी विखंडन}। एक घंटे बाद अधिकांश जुड़ चुके
होते हैं: \qty{30}{min} के आसपास के अर्ध-आयु काल पर लगभग $70/4 \approx 18$
केंद्र। छह घंटे बाद लगभग सब मरम्मत हो चुके होते हैं ($70/2^{12} < 1$),
और केवल कुछ टिकाऊ, जटिल विखंडन बचते हैं।
\end{solution}

\begin{solution}{exo:b3:dna-repair:5}
(क) $10^{-5}\times 10^{-2}\times 10^{-3} = 10^{-10}$; $6.4\times
10^{9}\times 10^{-10} = 0.64$ उत्परिवर्तन प्रति विभाजन। (ख) \omterm{prop:b3:dna-repair:fidelity}{असंगति
मरम्मत} के बिना $10^{-7}$; प्रति विभाजन $640$। (ग) शुद्धि-पाठन के बिना
$10^{-8}$; प्रति विभाजन $64$। असंगति-मरम्मत उत्परिवर्ती अधिक खतरनाक
है, सामान्य से हज़ार गुना।
\end{solution}

\begin{solution}{exo:b3:dna-repair:6}
$k = \ln 2/\qty{0.25}{h} = \qty{2.8}{h^{-1}}$; $D = 10^{4}/24 =
\qty{420}{h^{-1}}$; $L^{*} = 420/2.8 \approx 150$ अक्षारीय स्थल किसी भी
क्षण उपस्थित। दस गुना धीमी मरम्मत: $L^{*} \approx \num{1500}$।
द्विशाख हर स्थल से एक बार गुज़रता है और वहाँ उसी क्षण उपस्थित
\omterm{def:b3:dna-repair:lesions}{विक्षतियों} से मिलता है; जीनोम भर औसत लेने पर वह उनमें से लगभग $L^{*}$
से मिलता है --- सामान्यतः कोई $150$, औषध के साथ $\num{1500}$।
\end{solution}

\begin{solution}{exo:b3:dna-repair:7}
$k_{\text{NHEJ}} = \ln 2/0.5 = \qty{1.39}{h^{-1}}$, $k_{\text{HR}} =
\ln 2/4 = \qty{0.17}{h^{-1}}$; HR अंश $0.17/1.56 = 0.11$;
अर्ध-आयु $\ln 2/1.56 = \qty{27}{min}$। कोई ढहा द्विशाख
\emph{एक-सिरे वाला} विखंडन देता है: NHEJ के जोड़ने के लिए कोई दूसरा
सिरा ही नहीं है, इसलिए द्विशाख केवल सहोदरी से पुनर्संयोजन ही बहाल कर
सकता है --- और वैसे भी यही एकमात्र यथार्थ पथ है।
\end{solution}

\begin{solution}{exo:b3:dna-repair:8}
(क) लिंच: सूक्ष्म-उपग्रह अस्थिर (अर्बुद कोशिकाओं के बीच लंबाइयाँ भिन्न),
बृहदांत्र, गर्भाशयांतःकला, आमाशय और अंडाशय के कर्कट, धूप के प्रति
सामान्य अनुक्रिया। (ख) XP-A: सूक्ष्म-उपग्रह स्थिर, न्यूनतम संपर्क पर
अत्यधिक धूप-दाह, बचपन में त्वचा और नेत्र के कर्कट, और इस समूह में
उत्तरोत्तर तंत्रिकीय अपह्रास। (ग) XP-V: सूक्ष्म-उपग्रह स्थिर, उच्छेदन
मरम्मत सामान्य इसलिए धूप-दाह हल्का, पर जो द्विलक उच्छेदन से बच जाते
हैं उन्हें त्रुटि-प्रवण पॉलीमरेज़ लाँघती हैं --- त्वचा-कर्कट, समूह~A
से बाद में आरंभ होने वाले।
\end{solution}

\begin{solution}{exo:b3:dna-repair:9}
यकृत-कोशिकाएँ: Cre ने दोनों युग्मविकल्पियों से बहिर्वेशन~3 उच्छेदित कर
दिया है (एक वलय के रूप में, जो खो जाता है), इसलिए जीन जन्म से हर
यकृत-कोशिका में निष्क्रिय है --- एक यकृत-विशिष्ट निष्क्रियण।
तंत्रिका-कोशिकाएँ: कोई Cre नहीं, फ़्लॉक्स्ड युग्मविकल्पी अक्षत और
प्रकार्यशील। उलटी दिशाएँ: Cre बहिर्वेशन को प्रतिलोमित करता है और फिर
वापस, अनंत काल तक, इसलिए किसी भी समय लगभग आधी कोशिकाएँ प्रतिलोमित,
अप्रकार्य बहिर्वेशन ढोती हैं --- एक विचित्रता, कोई साफ़ विलोपन नहीं।
\end{solution}

\begin{solution}{exo:b3:dna-repair:10}
जैसे-जैसे LexA कटता है उसकी सांद्रता धीरे-धीरे गिरती है। जो प्रचालक
LexA से कमज़ोर बँधते हैं वे थोड़ी-सी गिरावट पर खाली हो जाते हैं ---
\emph{uvrA}, \emph{uvrB}, \emph{recA} --- इसलिए यथार्थ मरम्मत पहले
शुरू होती है; \emph{sulA} और त्रुटि-प्रवण पॉलीमरेज़ों के प्रचालक LexA
से कसकर बँधते हैं और तभी खाली होते हैं जब वह लगभग समाप्त हो जाए, यानी
जब क्षति लंबे समय तक बनी रही हो। जिस जीवाणु में त्रुटि-प्रवण
पॉलीमरेज़ न हों वह उन \omterm{def:b3:dna-repair:lesions}{विक्षतियों} से रुके द्विशाखों पर मर जाएगा जिन्हें
वह लाँघ नहीं सकता, और, बिना किसी प्रेरित उत्परिवर्तनकारिता के,
प्रतिजैविक के प्रति प्रतिरोध कहीं धीमे विकसित करेगा --- जीवाणु के लिए
घाटा, रोगी के लिए लाभ; इसी कारण \omterm{prop:b3:dna-repair:sos}{SOS अनुक्रिया} के निरोधकों पर अध्ययन हो रहा है।
\end{solution}

\begin{solution}{exo:b3:dna-repair:11}
शृंखला: PARP निरुद्ध $\to$ एकलड़ी विखंडन बने रहते हैं $\to$ कोई
प्रतिकृति द्विशाख हर एक को एक-सिरे वाले \omterm{def:b3:dna-repair:lesions}{द्विलड़ी विखंडन} में बदल देता है
$\to$ मरम्मत के लिए \omterm{def:b3:dna-repair:dsb}{समजात पुनर्संयोजन} चाहिए $\to$ \omterm{def:b3:dna-repair:dsb}{BRCA} से वंचित अर्बुद
कोशिका वह नहीं कर सकती $\to$ गलत जुड़ाव, गुणसूत्र-ह्रास, मृत्यु।
सामान्य कोशिकाएँ, जिनमें एक प्रकार्यशील युग्मविकल्पी है, पुनर्संयोजन
करके बच जाती हैं। प्रतिरोध: (1) \emph{\omterm{def:b3:dna-repair:dsb}{BRCA}} में ऐसा दूसरा उत्परिवर्तन
जो वाचन-ढाँचा और कोई काम करता प्रोटीन बहाल कर दे --- अर्बुद का या रक्त
में घूमते अर्बुद DNA का अनुक्रमण करके, और विकिरण के बाद \omterm{def:b3:dna-repair:dsb}{Rad51} केंद्रों
की वापसी से पकड़िए; (2) सिरा-अपच्छेदन रोकने वाले प्रोटीनों का ह्रास
(53BP1--शील्डिन), जो BRCA1 के बिना ही पुनर्संयोजन अंशतः बहाल कर देता
है --- अभिरंजन पर उनकी अनुपस्थिति से और बहाल \omterm{def:b3:dna-repair:dsb}{Rad51} केंद्रों से पकड़िए;
साथ ही PARP1 के ऐसे उत्परिवर्तन जो उसे DNA पर फँसने से रोकें, या
औषध-बहिःस्रावी पंप, जिन्हें क्रमशः अनुक्रमण और अभिव्यक्ति से पकड़ा जाता है।
\end{solution}

\begin{solution}{exo:b3:dna-repair:12}
दो समजात द्विकुंडलियाँ, एक $A$ ढोती और दूसरी $a$, चिह्नक के आर-पार
एकल लड़ियाँ बदल लेती हैं: अब उस खंड पर हर एक के पास हर जनक से एक लड़ी
है --- $A/a$ असंगति वाला एक विषमयुग्मक। यदि \omterm{prop:b3:dna-repair:fidelity}{असंगति मरम्मत} प्रविष्ट
द्विकुंडली पर असंगति को $A$ में बदल दे, तो तीन क्रोमैटिड $A$ ढोते हैं
और एक $a$: $3{:}1$; $a$ में बदले तो $1{:}3$; बिना मरम्मत छोड़ी जाए तो
क्रोमैटिड अर्धसूत्री विभाजन के बाद की पहली समसूत्री बारी में $A$ और
$a$ पृथक करता है और दो भिन्न पुत्री बीजाणु देता है (अर्धसूत्रोत्तर
पृथक्करण, आठ-बीजाणु वाली ऐस्कस में $5{:}3$)। विषमयुग्मक खंड आरंभकारी
विखंडन के चारों ओर केवल कुछ सौ से कुछ हज़ार क्षार युग्म का होता है,
इसलिए केवल उसी के भीतर के चिह्नक परिवर्तित हो सकते हैं, और वे, रचना से
ही, जीन-विनिमय के बगल में हैं।
\end{solution}

\begin{solution}{pb:b3:dna-repair:1}
\textbf{1.} प्रति दिन $10^{4} + 400 + 10^{4} \approx 2\times 10^{4}$
\omterm{def:b3:dna-repair:lesions}{विक्षतियाँ}, प्रति वर्ष $7\times 10^{6}$ --- मोटे तौर पर जीनोम के प्रति
किलोक्षार प्रति वर्ष एक।
\textbf{2.} प्रति bp $10^{-5}\times 10^{-2}\times 10^{-3} = 10^{-10}$;
प्रति विभाजन $6.4\times 10^{9}\times 10^{-10} = 0.64$ उत्परिवर्तन।
\textbf{3.} प्रति विभाजन $0.64\times 0.015 \approx 0.01$ कूटलेखी
उत्परिवर्तन; $50$ विभाजनों में लगभग $0.5$।
\textbf{4.} प्रति bp $10^{-7}$, प्रति विभाजन $640$ उत्परिवर्तन। फिसली
कोशिकाएँ: मरम्मत के बिना $10^{-4}\times 40\times 10^{9} = 4\times 10^{6}$;
उसके साथ $10^{-7}\times 40\times 10^{9} = 4000$ --- हज़ार गुना
अंतर, जो \omterm{def:b3:dna-repair:mmr}{सूक्ष्म-उपग्रह अस्थिरता} के रूप में दिखता है।
\textbf{5.} यूरैसिल DNA के लिए पराया है, यूरैसिल-DNA \omterm{def:b3:dna-repair:excision}{ग्लाइकोसिलेज़} उसे
पहचानकर उच्छेदित कर देती है: मरम्मत हो जाती है। 5-मेथिलसाइटोसीन से बना
थाइमीन किसी G के सामने पड़ा एक सामान्य क्षार है; प्रतिकृति के बाहर उसे
गलत बताने वाला कुछ नहीं, और वह C$\to$T उत्परिवर्तन बन जाता है --- वही
क्रियाविधि जिसने जीनोम से \omterm{def:b3:chromatin-epigenetics:methylation}{CpG} घटाए हैं।
\textbf{6.} 8-ऑक्सोG \omterm{def:b3:dna-repair:excision}{ग्लाइकोसिलेज़} ऑक्सीकृत क्षार को तब हटाती है जब वह
अब भी C के सामने है (प्रतिकृति से पहले); दूसरी \omterm{def:b3:dna-repair:excision}{ग्लाइकोसिलेज़} 8-ऑक्सोG
के सामने गलत बैठे A को हटाती है (प्रतिकृति के बाद, G की मरम्मत का एक
और अवसर देकर); हाइड्रोलेज़ ऑक्सीकृत dGTP को नष्ट कर देती है ताकि वह A
के सामने न बैठे। जब तीनों विफल हों, 8-ऑक्सोG A से युग्मित होता है और
अगली बारी G$\cdot$C$\to$T$\cdot$A पारक्रमण जमा देती है।
\textbf{7.} सामान्य $k = \ln 2/2 = \qty{0.35}{h^{-1}}$; XP $k = \ln 2/70
= \qty{0.0099}{h^{-1}}$।
\textbf{8.} सामान्य $10^{5} e^{-2.77} \approx 6200$; XP $10^{5}
e^{-0.079} \approx \num{92000}$।
\textbf{9.} प्रति कोशिका $6.2$ और $92$ उत्परिवर्तन; कूटलेखी $0.09$ और $1.4$।
\textbf{10.} $500$ संपर्क: $47$ कूटलेखी उत्परिवर्तन (सामान्य) और $690$
(XP), प्रतिकृति से आए $0.5$ के सामने --- किसी सामान्य व्यक्ति में भी
त्वचा के उत्परिवर्तन-भार पर धूप का ही बोलबाला है।
\textbf{11.} कूटलेखी अनुक्रम $9.6\times 10^{7}$ bp; चालक जीन उसका
$4\times 10^{4}/9.6\times 10^{7} = 4.2\times 10^{-4}$ हैं। माध्य
चोटें $690\times 4.2\times 10^{-4} = 0.29$; $P(\ge 1) = 1 - e^{-0.29} =
0.25$। $10^{9}$ कोशिकाओं में से $2.5\times 10^{8}$ कोई चालक उत्परिवर्तन
ढोती हैं (सामान्य बच्चा: माध्य $0.02$, \qty{2}{\%})।
\textbf{12.} कोई द्विलक उत्परिवर्तन नहीं है; वह तभी बनता है जब कोई
द्विशाख उसे त्रुटि-प्रवण लाँघ से नकल करे, इसलिए अनेक अमरम्मत द्विलकों
वाली विभाजनशील कोशिकाएँ ही उत्परिवर्तित होती हैं। पराबैंगनी प्रकाश
केवल त्वचा और नेत्र तक पहुँचता है; भीतरी ऊतकों को कोई द्विलक नहीं
मिलते और XP में वे लगभग सामान्य रहते हैं।
\textbf{13.} $70$ विखंडन; केवल NHEJ के साथ \qty{1}{h} बाद $70\times
2^{-2} \approx 18$ केंद्र।
\textbf{14.} $k_{\text{NHEJ}} = \qty{1.39}{h^{-1}}$, $k_{\text{HR}} =
\qty{0.17}{h^{-1}}$; HR अंश $0.17/1.56 = 0.11$।
\textbf{15.} अर्ध-आयु $\ln 2/1.56 = \qty{27}{min}$; \qty{1}{h} बाद
$70\,e^{-1.56} \approx 15$ विखंडन।
\textbf{16.} प्रति कोशिका $70\times 0.015\times 0.7 \approx 0.7$ जीन
बिगड़ते हैं।
\textbf{17.} $n = 70$: $2415$ युग्म $\times 3\times 10^{-4} = 0.72$
स्थानांतरण। $n = 7$: $21\times 3\times 10^{-4} = 0.006$। विखंडनों की
संख्या $\propto D$ है पर युगपत् विखंडनों के \emph{युग्मों} की संख्या
$\propto n^{2} \propto D^{2}$ है।
\textbf{18.} $D = \qty{7}{h^{-1}}$, $L^{*} = 7/1.39 \approx 5$ विखंडन
एक समय पर उपस्थित: $10$ युग्म, $0.003$ स्थानांतरण --- उसी मात्रा को
एक झटके में देने से दो सौ गुना कम। खंडीकरण सामान्य ऊतक को खंडों के बीच
मरम्मत करने देता है और युगपत् विखंडन थोड़े रखता है।
\textbf{19.} अक्षत मरम्मत: \qty{2}{Gy} पर $e^{-2/1.5} = 0.26$,
\qty{6}{Gy} पर $e^{-4} = 0.018$। NHEJ नहीं: $e^{-4} = 0.018$ और
$e^{-12} = 6\times 10^{-6}$। $D_{0}$ वह मात्रा है जो प्रति कोशिका औसतन
एक घातक अमरम्मत \omterm{def:b3:dna-repair:lesions}{विक्षति} छोड़ती है (उत्तरजीविता $1/e$)।
\textbf{20.} दो-सिरे वाले G2 विखंडनों का वह \qty{11}{\%} जिसकी मरम्मत
HR यथार्थ ढंग से करता, अब NHEJ से जाता है, और हर एक-सिरे वाला
द्विशाख-विखंडन --- जिसकी सही मरम्मत NHEJ कर ही नहीं सकता --- गलत जुड़
जाता है: वर्षों में जीनोम विलोपन, स्थानांतरण और प्रति-संख्या के बदलाव
जमा करता है, यानी \emph{\omterm{def:b3:dna-repair:dsb}{BRCA}} अर्बुदों के अभिलाक्षणिक घाव।
\textbf{21.} प्रति दिन $10^{4}\times 0.01 = 100$ एक-सिरे वाले विखंडन।
एक-सिरे वाले विखंडन का कोई साथी सिरा नहीं होता; NHEJ या तो उसे छोड़
सकता है या किसी असंबंधित सिरे से जोड़ सकता है, जिससे कोई स्थानांतरण या
कोई ह्रास बनता है।
\textbf{22.} एक प्रकार्यशील \emph{BRCA2} युग्मविकल्पी पुनर्संयोजन के
लिए पर्याप्त प्रोटीन बना देता है: सामान्य कोशिकाएँ अपने द्विशाख-विखंडनों
की मरम्मत करके निरोधक से बच जाती हैं।
\textbf{23.} पहले के नीचे की ओर एक दूसरा ढाँचा-विस्थापन वाचन-ढाँचा बहाल
कर देता है और छोटा पर अंशतः प्रकार्यशील BRCA2 देता है, इसलिए
पुनर्संयोजन लौट आता है और औषध अब नहीं मारती। प्रतिरोधी अर्बुद में या
घूमते अर्बुद DNA में \emph{BRCA2} का अनुक्रमण करके, और प्रकार्यात्मक
रूप से विकिरण के बाद \omterm{def:b3:dna-repair:dsb}{Rad51} केंद्रों के फिर से प्रकट होने से पकड़िए।
\textbf{24.} मारने के लिए \omterm{prop:b3:dna-repair:fidelity}{असंगति मरम्मत} को क्षारकृत क्षारों में उलझना
पड़ता है; उससे वंचित कोई लिंच अर्बुद सहिष्णु है --- औषध उसे मारने में
विफल रहती है (ऐसे अर्बुदों का उपचार दूसरे उपायों से किया जाता है,
प्रतिरक्षा-चिकित्सा सहित, जिसके प्रति उनका भारी उत्परिवर्तन-भार उन्हें
संवेदी बनाता है)।
\textbf{25.} XP कोशिका में द्विशाख पर लगभग \num{92000} द्विलक; प्रति
क्षार युग्म प्रति विभाजन $10^{-10}$ उत्परिवर्तन; HR G2 विखंडनों का
\qty{11}{\%} मरम्मत करता है।
\end{solution}
